% -*- coding: utf-8 -*-
% !TEX program = xelatex
\documentclass[plain,noanswer,medmath]{../../template/jnuexam}
\usepackage{../../template/kysx}

\begin{document}


\examtitle{1993}{5}


\loadproblems{4}{1993P4}%载入试卷四


\makepart{填空题}{本题共5小题，每小题3分，满分15分}


\begin{problem}
$\lim_{n \to \infty} \left[\sqrt{1 + 2 + \cdots + n} - \sqrt{1 + 2 + \cdots + (n - 1)} \right] =$ \fillin{}．
\end{problem}


%【类似试卷四第一(2)题】
\begin{problem}
已知$y = f\left(\frac{3x - 2}{3x + 2} \right)$，$f'(x) = \arcsin x^2$，
则$\left. \frac{\dy}{\dx} \right|_{x = 0} =$ \fillin{}．
\end{problem}


\begin{problem}
$\int \frac{\dx}{(2 - x)\sqrt{1 - x}} =$ \fillin{}．
\end{problem}


\useproblem{4}{1}{4}%【同试卷四第一(4)题】


\begin{problem}
设$10$件产品有$4$件不合格品，从中任取两件，已知所取两件产品中有一件是不合格品，
则另一件也是不合格品的概率为 \fillin{}．
\end{problem}


\makepart{选择题}{本题共5小题，每小题3分，满分15分}


\useproblem{4}{2}{1}%【同试卷四第二(1)题】


\useproblem{4}{2}{2}%【同试卷四第二(2)题】


\begin{problem}
若$\alpha_1$, $\alpha_2$, $\alpha_3$, $\beta_1$, $\beta_2$都是$4$维列向量，
且$4$阶行列式
$$|\alpha_1, \alpha_2, \alpha_3, \beta_1| = m,\qquad
  |\alpha_1, \alpha_2, \beta_2, \alpha_3| = n,$$
则$4$阶行列式$|\alpha_3,\alpha_2,\alpha_1,\beta_1 + \beta_2|$等于\pickout{}
\begin{abcd}
\item $m + n$．      
\item $- (m + n)$．      
\item $n - m$．      
\item $m - n$．
\end{abcd}
\end{problem}


\begin{problem}
设$\lambda = 2$是非奇异矩阵$A$的一个特征值，
则矩阵$\left(\tfrac{1}{3}A^2\right)^{-1}$有一特征值等于\pickout{}
\begin{abcd}
\item $\frac{4}{3}$．			
\item $\frac{3}{4}$．			
\item $\frac{1}{2}$．			
\item $\frac{1}{4}$．
\end{abcd}
\end{problem}


\begin{problem}
设随机变量$X$与$Y$均服从正态分布，$X\sim N(\mu ,4^2)$，$Y\sim N(\mu ,5^2)$，记
$$p_1 = P\{X \le \mu - 4\},\qquad p_2 = P\{Y \ge \mu + 5\},$$
则\pickout{}
\begin{abcd}
\item 对任何实数$\mu$，都有$p_1 = p_2$．		
\item 对任何实数$\mu$，都有$p_1 < p_2$．
\item 只对$\mu$的个别值，才有$p_1 = p_2$．	
\item 对任何实数$\mu$，都有$p_1 > p_2$．
\end{abcd}
\end{problem}


\makepart{}{本题满分5分}%【同试卷四第三题】

\useproblem{4}{3}{0}


\makepart{}{本题满分7分}%【同试卷四第四题】

\useproblem{4}{4}{0}


\makepart{}{本题满分7分}

\begin{problem}
已知某厂生产$x$件产品的成本为$C = 25000 + 200x + \frac{1}{40}x^2$（元）．问：
\begin{enumerate}
\item 要使平均成本最小，应生产多少件产品？
\item 若产品以每件$500$元售出，要使利润最大，应生产多少件产品？
\end{enumerate}
\end{problem}


\makepart{}{本题满分6分}

\begin{problem}
设$p$,$q$是大于$1$的常数，且$\frac{1}{p} + \frac{1}{q} = 1$．
证明：对于任意$x > 0$，有$\frac{1}{p} x^p + \frac{1}{q} \ge x$．
\end{problem}


\makepart{}{本题满分13分}

\begin{problem}
运用导数的知识作函数$y = (x + 6)\e^{\frac{1}{x}}$的图形．
\end{problem}


\makepart{}{本题满分8分}

\begin{problem}
已知$3$阶矩阵$A$的逆矩阵为$A^{-1} = \begin{pmatrix}
 1&1&3 \\
 1&2&1 \\
 1&1&1
\end{pmatrix}$，试求其伴随矩阵$A^*$的逆矩阵．
\end{problem}


\makepart{}{本题满分8分}

\begin{problem}
设$A$是$m \times n$矩阵，$B$是$n \times m$矩阵，$E$是$n$阶单位矩阵（$m > n$）．
已知$BA = E$，试判断$A$的列向量组是否线性相关？为什么？
\end{problem}


\makepart{}{本题满分8分}

\begin{problem}
设随机变量$X$和$Y$独立，都在区间$[1,3]$上服从均匀分布；
引进事件$A = \{X \le a\}$，$B = \{Y > a\}$．\par
\begin{enumerate*}
\item 已知$P(A \cup B) = \frac{7}{9}$，求常数$a$；
\item 求$\frac{1}{X}$的数学期望．
\end{enumerate*}
\end{problem}


\makepart{}{本题满分8分}%【同试卷四第十一题】

\useproblem{4}{11}{0}


\end{document}
